\documentclass[12pt]{article}
\usepackage{amsmath, amssymb}
\usepackage[margin=1in]{geometry}
\setlength{\parskip}{6pt}
\title{QUBO Formulation: Weighted Set Packing Problem}
\date{}
\begin{document}
\maketitle

\subsection*{Weighted Set Packing Problem}

\paragraph{Problem Statement.}
Let $U$ be a finite universe of elements and let $\mathcal{S} = \{S_0, S_1, \dots, S_{m-1}\}$ be a family of subsets of $U$, where each subset $S_j \subseteq U$ is associated with a non-negative weight $w_j \ge 0$. The Weighted Set Packing Problem seeks to select a subfamily of pairwise disjoint subsets $\mathcal{S}' \subseteq \mathcal{S}$ such that the total weight of the selected subsets is maximized. Formally, the goal is to find a collection of indices $J \subseteq \{0, 1, \dots, m-1\}$ satisfying $S_j \cap S_k = \emptyset$ for all distinct $j, k \in J$, that maximizes $\sum_{j \in J} w_j$.

\paragraph{Decision Variables.}
For each subset index $j \in \{0, 1, \dots, m-1\}$, introduce a binary decision variable $x_j \in \{0, 1\}$. The variable $x_j$ takes the value $1$ if subset $S_j$ is selected, and $0$ otherwise. The vector $x = (x_0, x_1, \dots, x_{m-1}) \in \{0, 1\}^m$ encodes a candidate solution.

\paragraph{QUBO Derivation.}
\textbf{Step 1 --- Objective.} The original problem is a maximization:
\begin{align}
    \max_{x \in \{0,1\}^m} \sum_{j=0}^{m-1} w_j x_j.
\end{align}
To cast this as a QUBO minimization, we negate the objective:
\begin{align}
    \min_{x \in \{0,1\}^m} \left( -\sum_{j=0}^{m-1} w_j x_j \right).
\end{align}

\textbf{Step 2 --- Constraints.} The pairwise disjointness requirement is equivalent to enforcing that for every element $e \in U$, at most one selected subset contains $e$. This yields the following family of linear constraints:
\begin{align}
    \sum_{j : e \in S_j} x_j \le 1, \quad \forall e \in U.
\end{align}

\textbf{Step 3 --- Penalty Terms.} Since the variables $x_j$ are binary, the constraint $\sum_{j : e \in S_j} x_j \le 1$ is violated if and only if there exist distinct indices $j, k$ such that $e \in S_j$, $e \in S_k$, and $x_j = x_k = 1$. Thus, a violation for element $e$ is detected by the product $x_j x_k$ for any pair $j < k$ sharing element $e$. We introduce a penalty weight $\lambda > 0$ and construct the penalty function:
\begin{align}
    P(x) = \lambda \sum_{e \in U} \sum_{j < k : e \in S_j \cap S_k} x_j x_k.
\end{align}
Expanding the sum over elements, the coefficient of a specific product $x_j x_k$ (with $j < k$) is $\lambda$ multiplied by the number of elements shared by $S_j$ and $S_k$. Let $|S_j \cap S_k|$ denote the cardinality of the intersection. The penalty can be rewritten as a sum over all pairs:
\begin{align}
    P(x) = \lambda \sum_{0 \le j < k \le m-1} |S_j \cap S_k| \, x_j x_k.
\end{align}

\textbf{Step 4 --- QUBO Hamiltonian.} Combining the negated objective and the penalty term, the QUBO Hamiltonian $H(x)$ is:
\begin{align}
    H(x) = -\sum_{j=0}^{m-1} w_j x_j + \lambda \sum_{0 \le j < k \le m-1} |S_j \cap S_k| \, x_j x_k.
\end{align}

\textbf{Step 5 --- Q Matrix and Offset.} The Hamiltonian is expressed in the standard QUBO form $H(x) = \sum_{j \le k} Q_{j,k} x_j x_k + c$. Reading off the coefficients from the general expression above, we obtain:
\begin{align}
    Q_{j,j} &= -w_j, \quad \forall j \in \{0, \dots, m-1\}, \\
    Q_{j,k} &= \lambda \, |S_j \cap S_k|, \quad \forall 0 \le j < k \le m-1, \\
    c &= 0.
\end{align}

\paragraph{Penalty Weight Justification.}
The penalty weight $\lambda$ must be chosen sufficiently large to ensure that any constraint violation incurs a cost exceeding the maximum possible gain from the objective. Specifically, let $W_{\max} = \max_{j} w_j$. If a solution $x$ violates a constraint, there exists at least one pair $j, k$ with $x_j = x_k = 1$ and $S_j \cap S_k \neq \emptyset$. The penalty contribution is at least $\lambda$. By removing one of the conflicting variables from an infeasible solution, we obtain a feasible solution whose weight is reduced by at most $W_{\max}$. Therefore, choosing
\begin{align}
    \lambda > W_{\max}
\end{align}
ensures that the energy of any infeasible solution is strictly greater than that of a corresponding feasible solution, guaranteeing that the global minimum of $H(x)$ corresponds to a feasible solution.

\paragraph{Correctness Argument.}
Minimizing $H(x)$ over $\{0, 1\}^m$ recovers the optimal set packing solution. For any feasible solution $x$, the penalty term vanishes, so $H(x) = -\sum_{j} w_j x_j$, which is minimized exactly when the total weight is maximized. For any infeasible solution $x$, the penalty $P(x) \ge \lambda$. Since $\lambda > W_{\max}$, the energy increase due to the penalty outweighs any potential reduction in the objective value relative to a feasible solution. Consequently, the global minimum of $H(x)$ must occur at a feasible point, and among feasible points, the minimum energy corresponds to the maximum weight packing.

\paragraph{Illustrative Example.}
Consider a small instance with $m=3$ subsets and universe $U = \{a, b\}$. Define the subsets and weights as:
\begin{align}
    S_0 = \{a\}, \quad w_0 = 5; \qquad
    S_1 = \{a, b\}, \quad w_1 = 8; \qquad
    S_2 = \{b\}, \quad w_2 = 4.
\end{align}
The intersections are $S_0 \cap S_1 = \{a\}$, $S_1 \cap S_2 = \{b\}$, and $S_0 \cap S_2 = \emptyset$. Thus $|S_0 \cap S_1| = 1$, $|S_1 \cap S_2| = 1$, and $|S_0 \cap S_2| = 0$. Choosing $\lambda = 9$ (since $\max w_j = 8$), the Hamiltonian becomes:
\begin{align}
    H(x) = -5 x_0 - 8 x_1 - 4 x_2 + 9 x_0 x_1 + 9 x_1 x_2.
\end{align}
The Q matrix entries are:
\begin{align}
    Q_{0,0} &= -5, & Q_{0,1} &= 9, & Q_{0,2} &= 0, \\
    Q_{1,1} &= -8, & Q_{1,2} &= 9, & Q_{2,2} &= -4.
\end{align}
The feasible solutions are $x \in \{(0,0,0), (1,0,0), (0,1,0), (0,0,1), (1,0,1)\}$. Evaluating $H(x)$:
\begin{align}
    H(0,0,0) &= 0, & H(1,0,0) &= -5, & H(0,1,0) &= -8, \\
    H(0,0,1) &= -4, & H(1,0,1) &= -9.
\end{align}
The minimum energy is $-9$ at $x^* = (1, 0, 1)$, corresponding to selecting subsets $S_0$ and $S_2$ with total weight $5 + 4 = 9$. This is optimal since selecting $S_1$ alone yields weight $8$, which is less than $9$.
\end{document}